\chapter[Sermon 18. He Exhorts Them to Persevere in the True Faith, Propounding Most Ample Rewards.]{He Exhorts Them to Persevere in the True Faith, Propounding Most Ample Rewards.}
\label{sermon:018}
\markright{Sermon 18. He Exhorts Them to Persevere in the True Faith ...}

And they will know that I have loved you, because you have kept the words of my patience. Therefore, I will keep you from the hour of temptation that is coming upon the whole world to test those who dwell on the earth. Behold, I am coming shortly. Hold fast to what you have, so that no one takes away your crown.

A commendable virtue is praised in the congregation of Philadelphia: they have kept the word of Christ, not every word, but the word of Christ, and have not denied it. And he has begun to recount great rewards, which both he has given to this church, and is also ready to give to any other church similarly devoted to godly religion. For we are drawn by rewards.

First, he says, I will convert your enemies, so that they may become your friends and brethren, coming into the congregation to worship Christ, whom they have blasphemed until now, and submitting themselves humbly and lowly. As we read of Saint Paul, who in the 15th chapter of the first epistle to the Corinthians says that he is unworthy to be called an apostle, and so forth. And this is a wonderful benefit. For God is glorified by those who are converted, and the truth is proclaimed; lies and superstition are confounded. The saints cannot but be exceedingly glad. The faithful also are delivered out of the Devil’s clutches and are saved.

Then follows another benefit of God. God’s enemies will discover and find that the church, and every member of it, are God’s well-beloved children. God’s enemies suppose the faithful to be wicked, enemies of God, heretics, church robbers, hated by God, and unworthy to live. But they will understand that nothing is dearer to God than the church, for whom he gave his son, whom he also chose as his spouse, and has made a partaker of his kingdom.

But of this love of God, by which he, provoked by no merit of ours, but by his only grace and innate goodness, has joined himself to the church, all virtues proceed. That chiefly which immediately follows: that the church has kept the word of patience. John in his canonical epistle says, “not that we have loved God, but that he has loved us, and so forth.” Therefore, where the observance of the word of patience is joined as the cause of love, it must be religiously explained, that the favor of God and all our gifts are truly of grace, but yet that he, through that same grace, seems to requite and reward us for our efforts. The saints are not proud because of this, but humbly acknowledge and preach grace everywhere and in all things.

Again he commends the perseverance of the faithful in the true religion. You have kept, he says, the word of my patience. The word of patience is the Gospel of eternal salvation, which Paul also calls the word of the cross, and that for two reasons. First, because he describes the cross and patience of Christ by which we are saved. And again, he urges us also to bear the cross and patiently suffer with Christ, Matthew 16:24; 2 Timothy 2:3. Neither must anyone expect perseverance from one who is impatient. The Lord says in Luke 12:34, “In your patience you will possess your souls.” Therefore, either the pastor or the church of Philadelphia has kept the word of patience, namely, by retaining the patience of Christ in their hearts through faith, and by showing patience in words and sayings, and enduring much hardship in body. This indeed is the best way to keep churches and every faithful person safe and sound. Let them keep the word of Christ’s patience, and commit the rest to the Lord.

\index{Rome}For it follows: And I will keep you again from the hour of temptation, \&c. The hour of temptation is explained in two ways. Either he speaks of heresies and heretics, by whose talk, and crafty juggling, lewdness, and deceptiveness tempts the faith, simplicity, and integrity of the faithful. Of this, the Lord treats much in the 13th Chapter of Deuteronomy. Or else he speaks truly of the persecutions which the emperors of Rome have inflicted, among whom Trajan, a most mighty prince, issued severe proclamations against the Christians. Of this, Pliny also made mention in the 10th book of Epistles, the holder and one. But Christ preserved the church of Philadelphia, and keeps also at this day the faithful by his word and power in the perils of heretics and heresies, and finally of persecutions also: so that the faithful may stand sure in all controversies and receive nothing of heretics that is strange from God’s word, and also give no place in persecutions. Christ causes many times that the burden of persecution presses not so heavily. Therefore let us always be steadfast in God’s word, and permit the defense to our Lord God. He will not neglect us, \&c.

\index{Hebrews, Epistle to the}But because in temptations and afflictions the Lord often appears to our flesh to delay, and to seem to neglect us, we say that the Lord anticipates and adds, “Behold, I come shortly.” Shortly, I say, that is to say, in time—neither late nor hastily. We say this neither too soon nor too late, but in due time and season. If the Lord therefore seems to be slow, do not despair, for he will come at the opportune moment when he sees fit. Do not prescribe to him the manner and means of deliverance, but wait for the Lord’s appointed time. Read what wholesome things Paul has written concerning this matter at the end of the tenth chapter to the Hebrews, where a passage also from the second chapter of Habakkuk is cited.

And now he exhorts in few but most evident words to perseverance in piety, wherein they had excelled hitherto. And he says two things are to be held fast: that which thou hast. They had the gospel of Christ, and the word of eternal life, the true faith and godly religion. These things he commands to hold fast, and to persist in the religion one has received. And whilst he commanded them to keep that which they had, he signifies by the way that no new or other doctrine is to be looked for, but that this one received doth suffice. Let us not therefore think in the government of the church upon other laws, upon other traditions, than of the Gospel of Jesus Christ. This is sufficient for the church. After reasoning as it were of the loss, he says: Therefore must thou watch diligently and hold strongly the gospel, for if this is taken away thou art spoiled of thy crown. The crown is a token of virtue and victory.

\index{Ezekiel (Book of)}Rulers and those worthy of empire are crowned. The virgin loses her crown, that is, is defiled. Therefore, heretics, false prophets, and seducers take away the crown when they seduce and corrupt. Thus says the Lord: “You have gained honor and glory; see that no one takes it from you.” So we read that Paul spoke in 2 Colossians. Let no one take the victory from you. In Ezekiel 18, the Lord testifies that he will not impute righteousness to the righteous if he forsakes and abandons his righteousness. Let us therefore pray that we may always persevere in the word of the Lord.
